\documentclass[a4paper,12pt]{article}
\usepackage{color}
\usepackage{graphicx, gensymb}
\usepackage{amsfonts, cancel, amssymb}
\usepackage{amsmath, amsthm, array, esvect, esint}
\usepackage{typearea}
\usepackage[a4paper,left=2cm,right=2cm,top=2cm,bottom=3cm,bindingoffset=5mm]{geometry}
\newcommand\numberthis{\addtocounter{equation}{1}\tag{\theequation}}
\newcommand{\bra}{}
\newcommand{\kom}{}
\newcommand{\brak}{}
\newcommand{\braket}{}
\newcommand{\intx}{}
\newcommand{\intr}{}
\newcommand{\kla}{}
\newcommand{\klam}{}
\newcommand{\klammer}{}
\newcommand{\oper}{}
\newcommand{\tecxt}{}

\begin{document}

\title{06 Projektoren}
\author{Joachim Schwardt}
\date{\today}
\maketitle

\section{orthogonale Projektoren}
a)
\begin{proof}
$T$ ist beschränkt. Sei $\chi\in\mathcal{H}$:
$$\|T\chi\|=|\brak\langle {\phi}| {\chi}\rangle|\cdot\|\psi\|\le \|\phi\|\cdot\|\chi\|\cdot\|\psi\|=:M\|\chi\|$$
Also ist $\|T\|\le\|\phi\|\cdot\|\psi\|$. Es gilt zudem mit $\chi =\frac{\phi}{\|\phi\|}$, dass $\|T\|\ge\|\phi\|\cdot\|\psi\|$:
$$\|T\frac{\phi}{\|\phi\|}\|=\|\brak\langle {\phi}| {\frac{\phi}{\|\phi\|}}\rangle\cdot\psi\|=\|\phi\|\cdot\|\psi\|$$
Also gilt $\|T\| = \|\phi\|\cdot\|\psi\|$. Da $\phi,\psi\in\mathcal{H}\setminus\klammer\left\{ {0} \right\}$, gilt für den Kern:
\begin{align*}
T\chi &=\brak\langle {\phi}| {\chi}\rangle\psi=0\ \ \Rightarrow\ \ \ker T=\klammer\left\{ {\chi\in\mathcal{H}\ \vert\ \chi\bot\phi} \right\}=\klammer\left\{ {\phi} \right\}^\bot \\
T\chi &\propto\psi\ \ \Rightarrow\ \ \mathrm{im\,}T\subseteq\klammer\left\{ {\chi\in\mathcal{H}\ \vert\ \forall c\in\mathbb{C}:\chi=c\psi} \right\}= \mathrm{lin}\klammer\left\{ {\psi} \right\}\\
T\frac{\lambda\phi}{\|\phi\|^2}&=\lambda\psi\ \ \Rightarrow\ \ \mathrm{im\,}T = \mathrm{lin}\klammer\left\{ {\psi} \right\}
\end{align*}
Es gilt $T^\ast = \oper |\phi\rangle\langle \psi|$.
\begin{align*}
\brak\langle {\chi_1}| {T\chi_2}\rangle&=\brak\langle {\chi_1}| {\brak\langle {\phi}| {\chi_2}\rangle\psi}\rangle=\brak\langle {\psi}| {\chi_1}\rangle^\ast\brak\langle {\phi}| {\chi_2}\rangle=\brak\langle {\brak\langle {\psi}| {\chi_1}\rangle\phi}| {\chi_2}\rangle=\brak\langle {T^\ast\chi_1}| {\chi_2}\rangle
\end{align*}
\end{proof}
b)
\begin{proof}
$T$ ist normal, genau dann wenn $\phi = c\psi$ für ein $c\in\mathbb{C}$. 
\begin{align*}
T^\ast T&=\oper |\phi\rangle\langle \psi|\oper |\psi\rangle\langle \phi|=\|\psi\|^2\oper |\phi\rangle\langle \phi| \\
TT^\ast &=\oper |\psi\rangle\langle \phi|\oper |\phi\rangle\langle \psi|=\|\phi\|^2\oper |\psi\rangle\langle \psi|
\end{align*}
$T$ ist partiell isometrisch, genau dann wenn $\|\phi\|\cdot\|\psi\|=1$. Es gilt $(\ker T)^\bot=\klammer\left\{ {\phi} \right\}^{\bot\bot}=\mathrm{lin}\klammer\left\{ {\phi} \right\}$. Sei $T$ partiell isometrisch:
\begin{align*}
\|\phi\|&=\|T\phi\|=\|\brak\langle {\phi}| {\phi}\rangle\psi\|=\|\phi\|^2\|\psi\|
\end{align*}
Sei $\chi=\lambda\phi\in\mathrm{lin}\klammer\left\{ {\phi} \right\}$ und $\|\phi\|\cdot\|\psi\|=1$:
\begin{align*}
\|T\chi\|&=\|\brak\langle {\phi}| {\lambda\phi}\rangle\psi\|=|\lambda|\cdot\|\phi\|=\|\chi\|
\end{align*}
$T$ ist selbstadjungiert, genau dann wenn $\phi=\lambda\psi$ für ein $\lambda\in\mathbb{R}$.
\begin{align*}
T&=\oper |\psi\rangle\langle \phi|=\lambda\oper |\psi\rangle\langle \psi|=\oper |\lambda^\ast\psi\rangle\langle \psi|\overset{!}{=}\oper |\phi\rangle\langle \psi|=T^\ast
\end{align*}
$T$ ist eine Projektion, genau dann wenn $\brak\langle {\phi}| {\psi}\rangle=1$.
\begin{align*}
T^2&=T \ \ \Leftrightarrow\ \ \oper |\psi\rangle\langle \phi|\oper |\psi\rangle\langle \phi|=\oper |\psi\rangle\langle \phi|\ \ \Leftrightarrow\ \ \bra\langle {\phi\vert\psi  }\rangle = 1
\end{align*}
$T$ ist eine Orthoprojektion, genau dann wenn $\brak\langle {\phi}| {\psi}\rangle=\|\phi\|\cdot\|\psi\|=1$.
\end{proof}
\section{Beschränkte Operatoren}
a)
\begin{proof}
Sei $S:=T^\ast T$. Wähle für den anderen Fall $T\rightarrow T^\ast$. 
\begin{align*}
\brak\langle {\phi}| {S\phi}\rangle&=\brak\langle {T\phi}| {T\phi}\rangle=\|T\phi\|^2\ge 0 \\
S^\ast&=\kla\left( {T^\ast T} \right)^\ast=T^\ast T^{\ast\ast}=T^\ast T=S
\end{align*}
\end{proof}
b)
\begin{proof}
Setze $M:=\sup\klammer\left\{ {|\brak\langle {\phi}| {T\phi}\rangle| \ \vert\ \|\phi\|\le 1} \right\}$. Es folgt $M\le\|T\|$:
$$|\brak\langle {\phi}| {T\phi}\rangle| \le \|\phi\|\cdot\|T\phi\|\le\|T\|\cdot\|\phi\|^2$$ 
Für $\|\phi\|, \|\psi\|\le 1$ gilt:
\begin{align*}
\brak\langle {\phi+\psi}| {T(\phi+\psi)}\rangle-\brak\langle {\phi-\psi}| {T(\phi-\psi}\rangle&=2\brak\langle {\phi}| {T\psi}\rangle+2\brak\langle {\psi}| {T\phi}\rangle=4\mathrm{Re\,}\brak\langle {\phi}| {T\psi}\rangle
\end{align*}
Mit der Parallelogrammidentität folgt daraus:
\begin{align*}
4\mathrm{Re\,}\brak\langle {\phi}| {T\psi}\rangle &\le M\kla\left( {\|\phi+\psi\|^2+\|\phi-\psi\|^2} \right)=2M\kla\left( {\|\phi\|^2+\|\psi\|^2} \right)\le 4M
\end{align*}
Es gibt $\lambda\in\mathbb{C}$ mit $|\lambda|=1$ und $\lambda\brak\langle {\phi}| {T\psi}\rangle=|\brak\langle {\phi}| {T\psi}\rangle|$. Also folgt:
\begin{align*}
|\brak\langle {\phi}| {T\psi}\rangle|&=\mathrm{Re\,}|\brak\langle {\phi}| {T\psi}\rangle|=\mathrm{Re\,}\lambda\brak\langle {\phi}| {T\psi}\rangle=\mathrm{Re\,}\brak\langle {\lambda^\ast\phi}| {T\psi}\rangle\le M
\end{align*}
\end{proof}
c)
\begin{proof}
$(i)\Rightarrow(ii)$: \\
Es gilt:
\begin{align*}
\mathrm{im\,}T^\ast &\subseteq \overline{\mathrm{im\,}T^\ast}=\kla\left( {\mathrm{im\,}T^\ast} \right)^{\bot \bot}=\kla\left( {\ker T} \right)^\bot
\end{align*}
Damit folgt:
\begin{align*}
\brak\langle {\kla\left( {(TT^\ast)^2-TT^\ast} \right)\psi}| {\psi}\rangle&=\brak\langle {TT^\ast\psi}| {TT^\ast\psi}\rangle-\brak\langle {TT^\ast\psi}| {\psi}\rangle=\|TT^\ast\psi\|^2-\|T^\ast\psi\|^2=0
\end{align*}
$(ii)\Rightarrow(i)$: \\
Es gilt:
\begin{align*}
\|TT^\ast\psi\|^2&=\brak\langle {(TT^\ast)^2\psi}| {\psi}\rangle=\brak\langle {TT^\ast\psi}| {\psi}\rangle=\|T^\ast\psi\|^2
\end{align*}
$(ii)\Rightarrow(iii)$: \\
Es gilt: 
\begin{align*}
(T^\ast TT^\ast - T^\ast)^\ast (T^\ast TT^\ast - T^\ast)&=(TT^\ast T-T) (T^\ast TT^\ast - T^\ast) \\
&=(TT^\ast T)(T^\ast TT^\ast)-(TT^\ast T)T^\ast - T(T^\ast TT^\ast)+TT^\ast \\
&=(TT^\ast)^3-2(TT^\ast)^2+TT^\ast =0
\end{align*}
Wegen der $C^\ast -$Eigenschaft folgt:
\begin{align*}
\|T^\ast TT^\ast -T^\ast\|^2=\|(T^\ast TT^\ast -T^\ast)^\ast (T^\ast TT^\ast - T^\ast)=0\ \ \Rightarrow\ \ T^\ast = T^\ast TT^\ast
\end{align*}
$(iii)\Rightarrow(ii)$: \\
Es gilt:
\begin{align*}
(TT^\ast)^2&=TT^\ast TT^\ast = T(T^\ast TT^\ast)=TT^\ast
\end{align*}
$(iii)\Leftrightarrow(iv)$: \\
Es gilt:
\begin{align*}
T^\ast &=T^\ast TT^\ast\ \ \Leftrightarrow\ \ T=(T^\ast TT^\ast)^\ast = TT^\ast T\ \ \Leftrightarrow\ \ T^\ast\ \mathrm{partiell\ isometrisch}
\end{align*}
\end{proof}
d)
\begin{proof}
$(i)\Rightarrow(ii)$:\\
Aus der Polarisationsidentität folgt:
\begin{align*}
\brak\langle {T\phi}| {T\psi}\rangle&=\frac{1}{4}\sum_{k=1}^4i^k\|i^k T\phi+T\psi\|^2=\frac{1}{4}\sum_{k=1}^4i^k\|i^k \phi+\psi\|^2=\brak\langle {\phi}| {\psi}\rangle
\end{align*}
$(ii)\Rightarrow(i)$:\\
Es gilt:
\begin{align*}
\|T\phi\|^2&=\brak\langle {T\phi}| {T\phi}\rangle=\brak\langle {\phi}| {\phi}\rangle=\|\phi\|^2
\end{align*}
\end{proof}
e)
\begin{proof}
$(i)\Leftrightarrow(iv)$:\\
\begin{align*}
T\ \tecxt\text{normal}\ \ \Leftrightarrow\ \ T^\ast T = TT^\ast\ \ \Leftrightarrow\ \ T^\ast (T^\ast)^\ast = (T^\ast)^\ast T^\ast\ \ \Leftrightarrow\ \ T^\ast\ \tecxt\text{normal}
\end{align*}
$(i)\Rightarrow(ii)$:\\
Sei $\phi,\psi\in\mathcal{H}$:
\begin{align*}
\brak\langle {T\phi}| {T\psi}\rangle&=\brak\langle {\phi}| {T^\ast T\psi}\rangle=\brak\langle {\phi}| {TT^\ast\psi}\rangle=\brak\langle {T^\ast\phi}| {T^\ast\psi}\rangle
\end{align*}
$(ii)\Rightarrow(iii)$:\\
Sei $\phi\in\mathcal{H}$:
\begin{align*}
\|T\phi\|^2&=\brak\langle {T\phi}| {T\phi}\rangle=\brak\langle {T^\ast\phi}| {T^\ast\phi}\rangle=\|T^\ast\phi\|^2
\end{align*}
$(iii)\Rightarrow(i)$:\\
\begin{align*}
\brak\langle {(T^\ast T-TT^\ast)\phi}| {\phi}\rangle&=\brak\langle {T\phi}| {T\phi}\rangle-\brak\langle {T^\ast\phi}| {T^\ast\phi}\rangle=0
\end{align*}
\end{proof}
f)
\begin{proof}
$(i)\Leftrightarrow(iii)$:\\
\begin{align*}
T\ \tecxt\text{unitär}\ \ \Leftrightarrow\ \ T^{-1}=T^\ast\ \ \Leftrightarrow\ \ (T^{-1})^\ast=(T^\ast)^\ast\ \ \Leftrightarrow\ \ (T^\ast)^{-1}=(T^\ast)^\ast\ \ \Leftrightarrow\ \ T\ \tecxt\text{unitär}
\end{align*}
$(i)\Rightarrow(ii)$:\\
Da $T$ invertierbar ist, ist $T$ insbesondere surjektiv. Außerdem ist $T$ isometrisch:
\begin{align*}
\|T\phi\|^2&=\brak\langle {T\phi}| {T\phi}\rangle=\brak\langle {T^\ast T\phi}| {\phi}\rangle=\brak\langle {\phi}| {\phi}\rangle=\|\phi\|^2
\end{align*}
$(ii)\Rightarrow(i)$:\\
Da $T$ isometrisch ist, ist $T$ insbesondere injektiv. Mit der vorausgesetzten Surjektivität folgt also, dass $T$ bijektiv ist. Es gilt $T^\ast = T^{-1}$:
\begin{align*}
\brak\langle {(T^\ast T-I)\phi}| {\phi}\rangle&=\brak\langle {T\phi}| {T\phi}\rangle-\brak\langle {\phi}| {\phi}\rangle=\|T\phi\|^2-\|\phi\|^2=0
\end{align*}
\end{proof}
g)
\begin{proof}
$(i)\Rightarrow(ii)$:\\
Sei $\phi\in\mathcal{H  }$:
\begin{align*}
\brak\langle {T\phi}| {\phi}\rangle&=\brak\langle {\phi}| {T\phi}\rangle^\ast=\brak\langle {T^\ast\phi}| {\phi}\rangle^\ast=\brak\langle {T\phi}| {\phi}\rangle^\ast
\end{align*}
$(ii)\Rightarrow(i)$:\\
Sei $\phi\in\mathcal{H  }$:
\begin{align*}
\brak\langle {(T-T^\ast)\phi}| {\phi}\rangle&=\brak\langle {T\phi}| {\phi}\rangle-\brak\langle {\phi}| {T\phi}\rangle=\brak\langle {T\phi}| {\phi}\rangle-\brak\langle {T\phi}| {\phi}\rangle^\ast=0
\end{align*}
\end{proof}
\section{Projektionen}
a)
\begin{proof}
$(i)\Rightarrow(ii)$:\\
Betrachte:
\begin{align*}
\kla\left( {I-P} \right)^2&=I^2-IP-PI+P^2=I-2P+P=I-P
\end{align*}
$(ii)\Rightarrow(i)$:\\
Da $I-P$ eine Projektion ist, ist auch $P=I-(I-P)$ eine Projektion.
$(iii)\Rightarrow(i)$:\\
Sei $\phi,\psi\in\mathcal{H}$:
\begin{align*}
\brak\langle {P^2\psi}| {\phi}\rangle&=\brak\langle {\psi}| {(P^\ast)^2\phi}\rangle=\brak\langle {\psi}| {P^\ast\phi}\rangle=\brak\langle {P\psi}| {\phi}\rangle
\end{align*}
Alternativ für die Äquivalenz zwischen $(i)$ und $(iii)$:
\begin{align*}
(P^\ast)^2&=P^\ast P^\ast = (PP)^\ast = P^\ast
\end{align*}
\end{proof}
b)
\begin{proof}
Aus Aufgabe 3.4a) wissen wir bereits, dass $\ker P$ abgeschlossener Unterraum ist. Sei $\phi\in\mathcal{H}$. Es gilt $P\phi\in\ker (I-P)$, also ist im\,$P\subseteq\ker (I-P)$:
\begin{align*}
(I-P)P\phi&=IP\phi-P^2\phi=0
\end{align*}
Sei $\phi\in\ker (I-P)$. Dann ist $\phi=P\phi\in\tecxt\text{im\,}P$, also gilt $\ker(I-P)\subseteq\tecxt\text{im\,}P$:
\begin{align*}
0&=(I-P)\phi=\phi-P\phi
\end{align*}
Aus dem Übergang von $P\rightarrow I-P$ folgt dann mit a) die andere Behauptung.
\end{proof}
c)
\begin{proof}
Es gilt:
\begin{align*}
\|P\|&=\|P^2\|\le\|P\|^2\ \ \Rightarrow\ \ 0\le\|P\|(\|P\|-1)
\end{align*}
\end{proof}
d)
\begin{proof}
$(i)\Rightarrow(ii)$:\\
Es ist im\,$P\subseteq (\ker P)^\bot=(\tecxt\text{im\,}(I-P))^\bot$. Daraus folgt:
\begin{align*}
\brak\langle {P\phi}| {\phi}\rangle&=\brak\langle {P\phi}| {P\phi}\rangle+\brak\langle {P\phi}| {(I-P)\phi}\rangle=\|P\phi\|^2\ge 0
\end{align*}
$(ii)\Rightarrow(iii)$:\\
Positive Operatoren sind selbstadjungiert, vgl. Aufgabe 6.2g).\\
$(iii)\Rightarrow(iv)$:\\
Selbstadjungierte Operatoren sind normal.\\
$(iv)\Rightarrow(v)$:\\
Es gilt $(PP^\ast)^2=P(P^\ast P)P^\ast =P(P P^\ast)P^\ast =P^2(P^\ast)^2=PP^\ast$. Also ist $PP^\ast$ eine Projektion und somit $P$ eine partielle Isometrie. \\
$(v)\Rightarrow(vi)$:\\
Da $P$ partiell isometrisch ist, ist $P^\ast P$ eine Projektion:
\begin{align*}
(P^\ast P-P)^\ast (P^\ast P-P)&=(P^\ast P-P^\ast)(P^\ast P-P)=(P^\ast P)^2-(P^\ast)^2P-P^\ast P^2+P^\ast P=\\
&=(P^\ast P)^2-P^\ast P=0
\end{align*}
$(vi)\Rightarrow(vii)$:\\
Wegen der $C^\ast -$Eigenschaft gilt:
\begin{align*}
\|P\|^2&=\|P^\ast P\|=\|P\|\ \ \Rightarrow\ \ 0=\|P\|(\|P\|-1)
\end{align*}
$(vii)\Rightarrow(i)$:\\
Sei $\|P\|=1$, $\phi\in\ker P$, $\psi\in\tecxt\text{im\,}P$ und $\lambda\in\mathbb{C}$. Dann gilt:
\begin{align*}
P(\phi+\lambda\psi)&=P\phi+\lambda P\psi=\lambda\psi
\end{align*}
Damit folgt:
\begin{align*}
\|\lambda\psi\|^2&=\|P(\phi+\lambda\psi)\|^2\le \|\phi+\lambda\psi\|^2=\|\phi\|^2+2\tecxt\text{Re\,}\lambda\brak\langle {\phi}| {\psi}\rangle+\|\lambda\psi\|^2\\
&\Rightarrow\ \ \|\phi\|^2\ge -2|\lambda|\tecxt\text{Re\,}e^{i\alpha}|\brak\langle {\phi}| {\psi}\rangle|e^{i\beta}=-2|\lambda|\,|\brak\langle {\phi}| {\psi}\rangle|\cos(\alpha+\beta)
\end{align*}
Wähle also $|\lambda|=\mu >0$ beliebig und $\alpha=\pi-\beta$:
\begin{align*}
\|\phi\|^2\ge 2\mu |\brak\langle {\phi}| {\psi}\rangle|
\end{align*}
Damit diese Ungleichung für beliebig große $\mu$ gelten kann, muss $\brak\langle {\phi}| {\psi}\rangle=0$ gelten.
\end{proof}
\section{$C^\ast-$Algebra}
a)
\begin{proof}
Aus Aufgabe 1.2 wissen wir, dass $l^\infty$, $c$ und $c_0$ Banachräume sind. \\
Produktfolge zweier beschränkter Folgen ist beschränkt, also gilt $xy\in l^\infty$. Insbesondere ist $\|\cdot\|_\infty$ submultiplikativ. \\
Produktfolge zweier konvergenter Folgen ist konvergent, also ist $xy\in c$. \\
Produktfolge zweier Nullfolgen ist eine Nullfolge, also ist $xy\in c_0$. \\
Die Assoziativität, Bilinearität und Kommutativität der Multiplikation folgt direkt aus dem Assoziativ-, Distributiv. bzw. Kommutativgesetz im Körper $\mathbb{C}$. Das Gleiche gilt auch für die komplexe Konjugation. \\
Sei $x_k$ eine Folge in $\mathbb{C}$. Dann folgt die $C^\ast-$Eigenschaft:
\begin{align*}
|\overline{x_k}x_k|=||x_k|^2|=|x_k|^2
\end{align*}
$c_0$ hat kein Einselement, da $I\notin c_0 $
\end{proof}
b)
\begin{proof}
Sei $x_k$ eine Folge in $l^2$. Dann gilt $M_{ab}=M_aM_b$:
\begin{align*}
M_{ab}x&=(a_kb_kx_k)=M_a(b_kx_k)=M_aM_bx
\end{align*}
Weiterhin gilt $M_{a^\ast}=M_a^\ast$:
\begin{align*}
\brak\langle {y}| {M_ax}\rangle&=a_k\brak\langle {y_k}| {x_k}\rangle=\brak\langle {a_k^\ast y_k}| {x_k}\rangle=\brak\langle {M_{a^\ast}y}| {x}\rangle
\end{align*}
\end{proof}
\end{document}